\chapter*{Chapter 4}
\markboth{Chapter 4}{4 Implementation - Phase 2: Sprints 4-6}
\addcontentsline{toc}{chapter}{4 Implementation - Phase 2: Sprints 4-6}

\textsc{\textbf{\Huge Implementation - Phase 2}} \\ \\

\stepcounter{chapter}

% ============================================================
\section*{Introduction}
\addcontentsline{toc}{section}{Introduction}

This chapter presents the second phase of implementation work,
covering Sprints 4 through 6. While the first phase established
the core infrastructure for authentication, product catalog, and
data collection, this phase focused on building features that
directly serve the end users. The primary focus was on client
features like favorites and price alerts, merchant management
capabilities for offering products, and administrative tools
for overseeing the platform. Each sprint added significant
functionality that transformed the platform from a basic price
comparison tool into a fully-featured service with personalized
user experiences.

The work in these three sprints addressed the requirements
identified during the analysis phase and built upon the foundation
created in the first three sprints. The authentication system
from Sprint 1 was used to protect user-specific features. The
product catalog from Sprint 2 became the basis for personalization.
The scraping system from Sprint 3 fed the data that made price
alerts meaningful. This integration of previous work demonstrated
the value of building a solid foundation before adding layers
of functionality.

% ============================================================
\section*{Sprint 4: Client Personalization Features}
\addcontentsline{toc}{section}{Sprint 4: Client Personalization Features}

\subsection*{Sprint Overview}

Sprint 4 concentrated on features that make the platform useful
for returning clients. While visitors could search and compare
prices, authenticated clients gained the ability to save their
preferences, track products they were interested in, and receive
notifications when prices changed. These features increased user
engagement and provided additional value that encouraged users
to create accounts.

The three main client features implemented were favorites, price
alerts, and wishlists. Each served a different purpose in helping
clients manage their shopping experience. Favorites provided a
quick way to save interesting products for later reference. Price
alerts automated the process of monitoring price changes and
notifying users when products reached their target prices.
Wishlists allowed clients to organize products into custom
collections, useful for comparing multiple items or planning
purchases.

A critical requirement for these features was that they only
be available to authenticated clients. Visitors could browse and
compare prices, but saving favorites or creating alerts required
logging in. This distinction was enforced through the role-based
middleware already in place from Sprint 1, with the client role
being required for access to these endpoints.

The technology approach leveraged the existing database models
and Laravel relationships to implement these features efficiently.
Rather than creating entirely new data structures, the features
built upon the relationships already defined between products,
offers, and users. This approach maintained consistency across
the platform and reduced the complexity of the implementation.

Favorites worked as a simple bookmark system where clients could
mark any product as interesting with a single click. The system
stored the association between the client and the product,
allowing them to view all their favorites in one place. This
was particularly useful for clients who were comparing options
across multiple visits, as they could quickly return to products
they had previously identified as interesting.

The favorites feature utilized a many-to-many relationship
between clients and products. When a client added a product to
their favorites, the system created an entry in the favorites
table linking the two. This design allowed a client to have
any number of favorites and a product to be favorited by many
clients. The interface displayed the favorite status clearly,
allowing clients to easily add or remove items from their list.

Price alerts represented a more sophisticated feature that
combined user preferences with ongoing price monitoring. When a
client created an alert for a product, they specified a target
price they considered a good deal. The system would then monitor
the product's offers and notify the client when any offer fell
below their target price.

The alert system required background processing to check prices
regularly and trigger notifications. This was implemented using
scheduled tasks that ran periodically to evaluate active alerts.
When a price drop was detected, the system sent a notification
to the client through their preferred channel, typically SMS
since phone numbers were already verified during registration.

Wishlists provided organizational functionality for clients who
wanted to group products together. Unlike favorites which was a
single flat list, wishlists allowed clients to create multiple
named collections. This was useful for different purposes such
as organizing gifts for different occasions, comparing products
in a specific category, or tracking items across different
shopping trips.

The wishlist system included both the list containers and the
items within them. Clients could create any number of wishlists
with custom names. Each wishlist could contain multiple products.
The system also allowed clients to share their wishlists with
others through a unique link, useful for gift suggestions or
collaborative shopping.

\begin{table}[H]
\begin{center}
\begin{tabular}{|m{3cm}|m{6cm}|}
\hline
\textbf{Feature} & \textbf{Description} \\ \hline
Favorites & One-click saving of products to a personal list \\ \hline
Price Alerts & Automated monitoring with SMS notification on price drops \\ \hline
Wishlists & Custom collections of products for organization and sharing \\ \hline
\end{tabular}
\caption{Client Personalization Features}
\label{tab:client_features}
\end{center}
\end{table}

The controller layer for these features was organized around the
client domain. A FavoriteController handled adding and removing
favorites, retrieving the list of favorites, and checking whether
a specific product was already favorited. A PriceAlertController
managed the creation, modification, and deletion of alerts,
along with listing active alerts. A WishlistController handled
creating wishlists, adding and removing products from them, and
retrieving the contents.

Testing verified that all three features worked correctly.
Clients could add and remove favorites without errors. Price
alerts were created correctly and notifications were sent when
conditions were met. Wishlists could be created, populated, and
retrieved with all products properly associated. The access
controls worked as expected, preventing visitors from accessing
any of these features without authentication.

\begin{figure}[H]
\begin{center}
\fbox{Add screenshot: Favorites page / Price alerts / Wishlists}
\end{center}
\caption{Client Personalization Features}
\label{fig:client_features_interface}
\end{figure}

% ============================================================
\section*{Sprint 5: Merchant Management Features}
\addcontentsline{toc}{section}{Sprint 5: Merchant Management Features}

\subsection*{Sprint Overview}

Sprint 5 shifted focus to the merchant side of the platform.
While Sprint 4 served clients who wanted to find the best
prices, this sprint served the businesses that offered those
prices. Merchants needed tools to manage their product offerings,
subscribe to platform plans that provided access to different
feature levels, and track the performance of their listings.

The core merchant functionality centered on offer management.
Merchants could add new offers for products in the catalog,
update prices and availability for existing offers, and remove
offers that were no longer valid. This was distinct from the
scraping system in Sprint 3, which automatically collected offers
from external websites. Merchant offers were created manually
by the merchants themselves through their dashboard.

Subscription management provided the business model for the
platform. Merchants could subscribe to different plans that
granted access to various features. Basic plans might allow a
limited number of offers, while premium plans offered unlimited
listings with additional visibility features. The subscription
system tracked which plan each merchant was on and when their
current period expired.

The offer management system required careful integration with
the existing product catalog. Rather than creating new products,
merchants would add offers to existing products in the catalog.
This maintained the integrity of the canonical product data while
allowing merchants to contribute their pricing information.

When a merchant added a new offer, the system first checked
whether the referenced product existed in the catalog. If it
did, the offer was created with a link to that product. If the
product did not exist, the offer was held in a pending state
for administrator review. This approach prevented duplicate
products from proliferating while allowing legitimate new
products to be added through the merchant workflow.

The offer data included the price, URL to the product on the
merchant's website, availability status, and optional metadata
like stock quantity. The price was the critical field that
contributed to the price comparison functionality. The URL
enabled the redirect feature that sent clients to merchant
websites to complete purchases.

Subscription plans were defined with different feature sets and
pricing. The system tracked which plan each merchant had active,
along with the start and end dates of their current subscription
period. When a subscription was nearing expiration, the system
could send notifications to encourage renewal. The plan features
were enforced by checking the subscription status before allowing
certain actions like adding new offers.

Performance statistics gave merchants insight into how their
offers were performing. The system tracked metrics like how
many times each offer had been displayed in search results,
how many clicks had resulted in redirects to the merchant website,
and how many times the offer had been selected as the lowest
price. These metrics helped merchants understand the value
they were receiving from the platform.

\begin{table}[H]
\begin{center}
\begin{tabular}{|m{3cm}|m{6cm}|}
\hline
\textbf{Component} & \textbf{Function} \\ \hline
Offer Management & Add, update, and remove product offers \\ \hline
Subscription System & Plan-based access control and billing tracking \\ \hline
Performance Metrics & Click and view statistics for merchant offers \\ \hline
Merchant Profile & Business information and verification status \\ \hline
\end{tabular}
\caption{Merchant Feature Components}
\label{tab:merchant_components}
\end{center}
\end{table}

The merchant controller layer provided the API endpoints for
all merchant functionality. Endpoints existed for managing
offers, viewing and updating subscriptions, and retrieving
performance reports. The role-based access control ensured
that only users with the merchant role could access these
features, and merchants could only access their own data.

Testing confirmed that merchants could successfully create and
manage offers. The subscription system correctly enforced plan
limits. Performance metrics were accurately tracked and displayed.
The access controls properly restricted functionality to
authenticated merchants only.

\begin{figure}[H]
\begin{center}
\fbox{Add screenshot: Merchant dashboard / Offer management / Subscriptions}
\end{center}
\caption{Merchant Interface}
\label{fig:merchant_interface}
\end{figure}

% ============================================================
\section*{Sprint 6: Administrative Features}
\addcontentsline{toc}{section}{Sprint 6: Administrative Features}

\subsection*{Sprint Overview}

Sprint 6 completed the feature implementation with a focus on
administrative capabilities. While previous sprints built features
for clients and merchants, this sprint provided the tools needed
to manage and oversee the entire platform. Administrators needed
to validate products, manage users, configure the system, and
monitor overall platform health.

Product validation addressed a quality control need that emerged
from the scraping system in Sprint 3. When products were scraped
from merchant websites, they entered the system in a pending
state. Administrators reviewed these pending products and decided
whether to publish them to the public catalog. This human review
ensured that products met quality standards before becoming
visible to users.

User management gave administrators complete control over user
accounts. They could view all users in the system, update roles
to change user types, and deactivate accounts that violated
terms of service. This functionality was essential for maintaining
the integrity of the user community and handling problematic
accounts.

The platform dashboard provided a high-level view of system
health. Administrators could see metrics like total users,
active products, number of merchants, and scraping health.
This visibility helped identify problems quickly and track
the growth of the platform over time.

Category and brand management allowed administrators to maintain
the taxonomy that organized the product catalog. They could add
new categories, modify existing ones, and set parent-child
relationships for hierarchical organization. Similarly, brands
could be added, merged, or removed as the product catalog evolved.

The product validation workflow followed a clear sequence.
Scraped products first entered a pending queue. Administrators
could review each pending product, examining the name, description,
category, brand, and other attributes. They could then approve
the product to publish it, reject it to remove it from the queue,
or request changes if the data needed correction before approval.

The dashboard aggregated metrics from across the platform.
User counts showed how many clients, merchants, and other users
were registered. Product counts displayed how many products and
offers existed in the catalog. Scraping statistics showed how
many scripts were running and their recent success rates. These
metrics updated in near real-time to give administrators current
information.

The category management interface allowed hierarchical organization.
Categories could be created with parent categories to form trees.
For example, Electronics could contain Computers, which could
contain Laptops. This structure was used throughout the platform
for browsing and filtering products. The brand management was
simpler, maintaining a flat list of manufacturer names that could
be assigned to products.

\begin{table}[H]
\begin{center}
\begin{tabular}{|m{3.5cm}|m{5.5cm}|}
\hline
\textbf{Admin Function} & \textbf{Description} \\ \hline
Product Validation & Review and approve scraped products before publication \\ \hline
User Management & View, modify roles, and deactivate user accounts \\ \hline
Dashboard & Platform-wide metrics and health monitoring \\ \hline
Category Management & Hierarchical taxonomy for product organization \\ \hline
Brand Management & Manufacturer listing for product attribution \\ \hline
Scraping Oversight & Monitor and configure data collection scripts \\ \hline
\end{tabular}
\caption{Administrative Features Summary}
\label{tab:admin_features}
\end{center}
\end{table}

The administrative controller provided endpoints for all these
functions. Product validation endpoints allowed reviewing pending
items and making approval decisions. User management endpoints
supported searching, filtering, and modifying user records.
Dashboard endpoints aggregated and returned the various metrics.
Category and brand endpoints provided full CRUD operations.

Testing verified that administrators could access all their
designated features. Product validation workflow functioned
correctly from pending through approval or rejection. User
management operations worked including role changes and
deactivation. Dashboard metrics displayed accurate data. Category
and brand management maintained proper data integrity.

\begin{figure}[H]
\begin{center}
\fbox{Add screenshot: Admin dashboard / Product validation / User management}
\end{center}
\caption{Admin Interface}
\label{fig:admin_interface}
\end{figure}

% ============================================================
\section*{Conclusion}
\addcontentsline{toc}{section}{Conclusion}
% ============================================================

The second phase of implementation expanded the platform from
a basic comparison tool into a full-featured service. Sprint 4
gave clients personalized experiences through favorites, price
alerts, and wishlists. Sprint 5 empowered merchants to manage
their offerings through dedicated tools for products, subscriptions,
and analytics. Sprint 6 provided the administrative oversight
needed to maintain quality and platform health.

These three sprints built directly on the foundation established
in the first phase. The authentication system protected user-specific
features. The product catalog provided the data that personalization
operated on. The scraping system fed the products that merchants
and administrators managed. This integration demonstrated how
each phase contributed to the overall platform.

The platform was now functionally complete with all major
features implemented. Users could search and compare prices as
visitors, or create accounts to access personalized features.
Merchants could manage their product offerings and track performance.
Administrators could oversee the entire system. The remaining
work would focus on refinement, optimization, and preparing
for production deployment.